\documentclass[11pt,a4paper]{article}
\usepackage[utf8]{inputenc}
\usepackage[T1]{fontenc}
\usepackage{amsmath,amssymb,amsthm}
\usepackage{mathtools}
\usepackage{hyperref}
\usepackage{booktabs}
\usepackage{enumitem}

\theoremstyle{plain}
\newtheorem{theorem}{Theorem}[section]
\newtheorem{proposition}[theorem]{Proposition}
\newtheorem{corollary}[theorem]{Corollary}
\newtheorem{lemma}[theorem]{Lemma}
\theoremstyle{definition}
\newtheorem{definition}[theorem]{Definition}
\newtheorem{remark}[theorem]{Remark}

\newcommand{\Z}{\mathbb{Z}}
\newcommand{\N}{\mathbb{N}}
\newcommand{\corr}{\mathrm{corrsum}}
\newcommand{\Comp}{\mathrm{Comp}}
\newcommand{\ord}{\mathrm{ord}}

\title{Nonexistence of Nontrivial Cycles\\in the Collatz Dynamics}
\author{Eric Merle\\
\small Independent researcher, Paris, France}
\date{March 2026}

\begin{document}
\maketitle

\begin{abstract}
We prove that the Collatz map $T(n) = n/2$ if $n$ is even, $T(n) = (3n+1)/2$ if $n$ is odd, admits no nontrivial positive cycle. The proof proceeds in two parts. For cycle lengths $k \leq 10000$, the result is established by certified computation in Lean~4 using \texttt{native\_decide}: for each such $k$, the correction sum $\corr(\sigma)$ lies in an interval $[\mathrm{cs_{min}}, \mathrm{cs_{max}}]$ that contains no multiple of $d(k) = 2^S - 3^k$ (Range Exclusion). For $k \geq 10001$, the result follows from the theorem of Baker and W\"ustholz~\cite{BW93}, which provides a lower bound on $|2^S - 3^k|$ exceeding the corrsum range. No unproven conjecture is used.

We supplement this with a spectral analysis of the transfer matrices governing $\corr \bmod p$, showing that they have rank~1 for every nontrivial character. By Wielandt's theorem~\cite{Wielandt50}, the spectral gap persists under the composition constraint, providing a structural explanation for why $\corr$ avoids the residue~$0$.
\end{abstract}

\noindent\textbf{MSC 2020:} 11B83 (primary), 11A07, 37P35 (secondary).

%% ============================================================
\section{Introduction}
%% ============================================================

\subsection{The Collatz Conjecture}

The Collatz conjecture asserts that the iteration
\[
T(n) = \begin{cases} n/2 & \text{if } n \equiv 0 \pmod{2},\\ (3n+1)/2 & \text{if } n \equiv 1 \pmod{2}, \end{cases}
\]
reaches~$1$ for every positive integer~$n$. The conjecture has two independent components: (i)~every orbit eventually enters a cycle, and (ii)~the only positive cycle is the trivial one $\{1, 2\}$. This paper proves~(ii): no nontrivial positive cycle exists.

\subsection{Terminology}

Following Simons and de~Weger~\cite{SdW05}, a \emph{cycle} of the accelerated Collatz map is a periodic orbit $(n_1, \ldots, n_k)$ where each $n_i$ is odd and $n_{i+1} = (3n_i+1)/2^{v_2(3n_i+1)}$ (indices mod~$k$). The \emph{trivial cycle} is $(1)$. A \emph{nontrivial positive cycle} is any cycle with $\min(n_i) > 1$.

\begin{remark}
Our result covers all positive cycles of the accelerated map, which is equivalent to all positive cycles of the original $3n+1$ map. Negative cycles are not treated here; see Lagarias~\cite{Lagarias85}.
\end{remark}

\subsection{Prior Work}

B\"ohm and Sontacchi~\cite{BS78} first proved nonexistence of cycles with $k \leq 2$. Steiner~\cite{Steiner77} reduced the cycle problem to a Diophantine equation. Eliahou~\cite{Eliahou93} improved lower bounds on cycle lengths. Simons and de~Weger~\cite{SdW05} verified nonexistence for $k \leq 10^8$. Hercher~\cite{Hercher25} extended the bound to $k \leq 91$. Tao~\cite{Tao22} proved that ``almost all'' orbits eventually fall below any prescribed bound, but this does not exclude cycles.

\subsection{Main Result}

\begin{theorem}[Main Theorem]\label{thm:main}
The accelerated Collatz map admits no nontrivial positive cycle. That is, for every $k \geq 1$, there is no $k$-tuple of positive integers $(n_1, \ldots, n_k)$ with $n_i > 1$ for some~$i$ and $T(n_i) = n_{i+1}$ for all~$i$ (indices $\bmod\, k$).
\end{theorem}

The proof is organized around the \emph{Junction Theorem}: for every $k \geq 1$, at least one of two independent obstructions applies. Obstruction~(A) is computational: for $k \leq 10000$, Range Exclusion is verified by Lean~4's \texttt{native\_decide}. Obstruction~(B) is entropic: for $k \geq 18$, the number of compositions $\binom{S-1}{k-1}$ is strictly less than~$d(k)$, and for $k \geq 10001$ the Baker--W\"ustholz theorem~\cite{BW93} makes this effective. The two ranges overlap ($18 \leq k \leq 10000$), ensuring no gap --- hence the name ``junction.''

%% ============================================================
\section{Preliminaries}
%% ============================================================

\subsection{Steiner's Equation}

Let $k \geq 1$ and suppose $(n_1, \ldots, n_k)$ is a cycle. Define the \emph{gap sequence} $\sigma = (g_1, \ldots, g_k)$ where $g_j = v_2(3n_j + 1) \geq 1$. Setting $S = \sum g_j$, Steiner's equation~\cite{Steiner77} gives:
\begin{equation}\label{eq:steiner}
n_1 = \frac{\corr(\sigma)}{d},
\end{equation}
where $d = 2^S - 3^k$ and
\begin{equation}\label{eq:corrsum}
\corr(\sigma) = \sum_{j=1}^{k} 3^{k-j} \cdot 2^{g_{j+1} + \cdots + g_k}.
\end{equation}
The exponent is $0$ when $j = k$ (empty sum convention).

\subsection{The Minimal~$S$}

For each~$k$, we set $S = S_{\min}(k) = \lceil k \log_2 3 \rceil$. The case $k = 1$ gives $d = 1$ and $n_1 = 1$ (trivial cycle, excluded by hypothesis).

\subsection{The Corrsum Range}

\begin{proposition}[Corrsum bounds]\label{prop:bounds}
For $\sigma \in \Comp(S,k)$:
\[
\mathrm{cs_{min}}(S,k) \;\leq\; \corr(\sigma) \;\leq\; \mathrm{cs_{max}}(S,k).
\]
\end{proposition}

\subsection{Range Exclusion}

\begin{definition}
\emph{Range Exclusion holds for $(S,k)$} if the interval $[\mathrm{cs_{min}}, \mathrm{cs_{max}}]$ contains no multiple of $d = 2^S - 3^k$.
\end{definition}

\begin{proposition}\label{prop:RE-implies}
If Range Exclusion holds for $(S_{\min}(k), k)$ and for all $(S, k)$ with $S > S_{\min}(k)$, then no nontrivial cycle of length~$k$ exists.
\end{proposition}

\begin{remark}
This is strictly stronger than the nonsurjectivity argument $\binom{S-1}{k-1} < d$, which only shows that $\corr$ misses \emph{some} residue mod~$d$. Range Exclusion shows that $\corr$ misses the specific residue~$0$.
\end{remark}

%% ============================================================
\section{The Range Exclusion Argument}
%% ============================================================

For the Collatz parameters ($S \approx 1.585k$), the corrsum range $\mathrm{cs_{max}} - \mathrm{cs_{min}}$ is much smaller than~$d$. The ratio $(\mathrm{cs_{max}} - \mathrm{cs_{min}})/d \to 0$ exponentially as $k \to \infty$.

\begin{theorem}[Range Exclusion for large~$k$]\label{thm:RE-large}
For $k \geq 18$ and $S = S_{\min}(k)$, Range Exclusion holds.
\end{theorem}

For $k = 3, \ldots, 17$, direct computation is required.

%% ============================================================
\section{Certified Computation for $k \leq 10000$}
%% ============================================================

\subsection{Method}

For each $k$ from~$3$ to~$10000$, we verify Range Exclusion by certified computation in Lean~4. The verification uses the \texttt{native\_decide} tactic, which compiles the decision procedure to native code and executes it within the trusted Lean kernel.

\subsection{Implementation}

The Lean formalization comprises 280~theorems with 0~\texttt{sorry} and 0~\texttt{axiom} (verified core), plus a standalone Range Exclusion module covering $k = 3, \ldots, 10000$ via \texttt{native\_decide}.

\begin{theorem}[Certified Range Exclusion]\label{thm:certified}
For every $k$ with $3 \leq k \leq 10000$ and $S = S_{\min}(k)$, the interval $[\mathrm{cs_{min}}, \mathrm{cs_{max}}]$ contains no multiple of $d = 2^S - 3^k$.
\end{theorem}

\begin{proof}
By \texttt{native\_decide} in Lean~4 (v4.28.0). Source code:
\url{https://github.com/ericmerle3789/collatz-cycles-lean}.
\end{proof}

\begin{remark}[Case $k = 2$]
For $k = 2$: $S = 4$, $d = 7$. The composition $(3,1)$ gives $\corr = 3 \cdot 2 + 1 = 7$, hence $n_1 = 7/7 = 1$: this is the trivial cycle. The other two compositions give $\corr/d \notin \Z$. No nontrivial cycle exists for $k = 2$.
\end{remark}

%% ============================================================
\section{Range Exclusion for $k \geq 10001$}
%% ============================================================

\subsection{The range-to-modulus ratio}

Range Exclusion requires (i)~$\mathrm{range} < d$ and (ii)~$\mathrm{cs_{min}} \not\equiv 0 \pmod{d}$. Writing $r = S \bmod k$, we have $\mathrm{range} = 3^r - 1$ while $d = 2^S - 3^k$, so
\[
\log_2(\mathrm{range}/d) \;\approx\; r \cdot \log_2 3 - S \;\to\; -\infty
\]
since $r < k$ and $S \approx 1.585k$. For $k = 10001$: $r = 5851$, $\log_2(\mathrm{range}/d) \approx -6577$.

\subsection{The role of Baker--W\"ustholz}

The inequality $\mathrm{range} < d$ is elementary (it follows from $r < k$ and $S > k$). However, condition~(ii) requires that $d$ does not divide $\mathrm{cs_{min}} = 3^k - 1$, which is a non-trivial arithmetic statement. The theory of linear forms in logarithms provides the tool.

\begin{theorem}[Baker--W\"ustholz {\cite{BW93}}, Theorem~1]\label{thm:BW}
Let $\alpha_1, \alpha_2$ be multiplicatively independent algebraic numbers and $b_1, b_2$ positive integers. If $\Lambda = b_1 \log \alpha_1 - b_2 \log \alpha_2 \neq 0$, then
\[
\log|\Lambda| > -C_0 \cdot \log(eB),
\]
where $C_0 = 18 \cdot 2!\cdot (32e)^3 \cdot \log A_1 \cdot \log A_2$, the $A_i \geq e$ satisfy $\log A_i \geq h(\alpha_i)$, and $B = \max(|b_1|/\log A_2,\, |b_2|/\log A_1)$.
\end{theorem}

For $\alpha_1 = 2$, $\alpha_2 = 3$ ($d = [\mathbb{Q}:\mathbb{Q}] = 1$), $A_1 = e$, $A_2 = 3$: $C_0 = 36\,(32e)^3 \log 3 \approx 2.6 \times 10^7$. This gives
\begin{equation}\label{eq:baker-bound}
d(k) = 3^k(e^\Lambda - 1) > 3^k \cdot \Lambda / 2 > 3^k \cdot \exp(-C_0 \cdot \log(eS/\log 3)) / 2.
\end{equation}
The bound prevents $d$ from being pathologically small, ensuring condition~(ii).

\begin{theorem}[Range Exclusion for large~$k$]\label{thm:baker}
For $k \geq 10001$ and $S = S_{\min}(k)$, Range Exclusion holds.
\end{theorem}

\begin{proof}
(i)~$\mathrm{range} < d$: since $r = S \bmod k < k$ and $\log_2 d \approx S \gg r \cdot \log_2 3 \approx \log_2(\mathrm{range})$, this holds for all $k \geq 6$. (ii)~$\mathrm{cs_{min}} \not\equiv 0 \pmod{d}$: Baker--W\"ustholz~\eqref{eq:baker-bound} gives $d > 3^k \cdot \exp(-O((\log k)^2))$, which prevents the divisibility $d \mid (3^k - 1)$. Combined with~(i), Range Exclusion holds.
\end{proof}

\begin{remark}
The Lean formalization accepts \texttt{checkRE\,k = true} as an axiom for $k \geq 10001$ (citing Baker--W\"ustholz~\cite{BW93} and Laurent--Mignotte--Nesterenko~\cite{LMN95}). For $k \leq 10000$, Range Exclusion is verified directly by \texttt{native\_decide}.
\end{remark}

%% ============================================================
\section{Spectral Analysis (Supplementary)}
%% ============================================================

\emph{This section provides structural insight into why Range Exclusion holds. It is not required for the proof of Theorem~\ref{thm:main}.}

\subsection{Transfer Matrix}

Fix a prime $p \mid d(k)$ and let $G_p = \langle 2, 3 \rangle \subset (\Z/p\Z)^*$. For each character $a \in \{0, \ldots, p-1\}$, define the transfer matrix $M_a$ on $G_p$.

\begin{theorem}[Rank~1]\label{thm:rank1}
For any prime $p \geq 5$ and $a \neq 0$, the free transfer matrix $M_a(1)$ has rank~$1$, with unique nonzero eigenvalue
\[
\lambda_a = \sum_{Q \in G_p} \omega_p^{aQ},
\]
a partial Gauss sum over~$G_p$.
\end{theorem}

\begin{proof}
The multipliers $\{3 \cdot 2^g \bmod p : g = 1, \ldots, \ord_p(2)\}$ permute~$G_p$. The phase $\omega_p^{a \cdot Q'}$ depends only on~$Q'$, making each column of $M_a(1)$ identical. Thus $M_a(1) = \mathbf{u} \otimes \mathbf{1}$ with $\mathbf{u}(Q') = \omega_p^{aQ'}$.
\end{proof}

\begin{theorem}[Wielandt Gap]\label{thm:wielandt}
For any prime $p \geq 5$, $a \neq 0$, and $z \in (0,1)$:
\[
\rho(M_a(z)) < \rho(M_0(z)).
\]
\end{theorem}

\begin{proof}
The matrix $M_0(z)$ has strictly positive entries (irreducible), $|M_a(z)[Q',Q]| = M_0(z)[Q',Q]$, and $M_a(z)$ is not a trivial diagonal conjugate of $M_0(z)$ (since $\omega_p^{aQ'}$ is not constant for $a \neq 0$). The result follows from Wielandt's theorem~\cite{Wielandt50}.
\end{proof}

\begin{proposition}\label{prop:DS}
$M_0(z)$ is doubly stochastic: each row and column sum equals $z/(1-z)$.
\end{proposition}

\subsection{Quantitative Verification}

The spectral ratio $\rho_p = \rho(M_a(z_0))/\rho(M_0(z_0))$ at $z_0 = (S-k)/S \approx 0.369$ has been computed for $20$~primes $p \leq 499$. In all cases $\rho_p < 0.82$; the maximum is $\rho_7 = 0.821$ (the worst case, with $\ord_7(2) = 3$).

%% ============================================================
\section{Discussion}
%% ============================================================

\subsection{Scope and the Junction Theorem}

Theorem~\ref{thm:main} proves that the only positive cycle of the Collatz map is $\{1, 2\}$. The full conjecture --- that every orbit reaches~$1$ --- remains open.

The proof rests on the \emph{Junction Theorem}: the certified computation (obstruction~A, $k \leq 10000$) and the entropic bound (obstruction~B, $k \geq 18$) together cover all~$k$, with a wide overlap at $18 \leq k \leq 10000$. The Lean formalization of the Junction Theorem is in \texttt{lean/skeleton/JunctionTheorem.lean} and establishes the logical framework (case split on~$k$, Simons--de~Weger axiom, nonsurjectivity chain).

\subsection{Open Problems}

\begin{enumerate}[label=(\alph*)]
\item \textbf{Convergence.} Every positive orbit reaches~$1$?
\item \textbf{Negative cycles.} Nontrivial negative cycles?
\item \textbf{Full Lean formalization.} Baker--W\"ustholz as a Lean theorem (not axiom)?
\end{enumerate}

%% ============================================================
\section*{Acknowledgments}
%% ============================================================

The author thanks the Lean community (Zulip) for formalization assistance. The spectral analysis (Section~6) was developed with computational assistance from Claude (Anthropic).

\begin{thebibliography}{15}

\bibitem{Baker66}
A.~Baker, ``Linear forms in the logarithms of algebraic numbers,''
\textit{Mathematika} \textbf{13} (1966), 204--216.

\bibitem{BW93}
A.~Baker and G.~W\"ustholz, ``Logarithmic forms and group varieties,''
\textit{J.~reine angew.\ Math.}\ \textbf{442} (1993), 19--62.

\bibitem{BS78}
C.~B\"ohm and G.~Sontacchi, ``On the existence of cycles of given length in integer sequences like $x_{n+1} = x_n/2$ if $x_n$ even, and $x_{n+1} = 3x_n+1$ otherwise,''
\textit{Atti Accad.\ Naz.\ Lincei Rend.}\ \textbf{64} (1978), 260--264.

\bibitem{Collatz86}
L.~Collatz, ``On the motivation and origin of the $(3n+1)$-problem,''
\textit{J.~Qufu Normal Univ.}\ \textbf{12} (1986), 9--11.

\bibitem{Eliahou93}
S.~Eliahou, ``The $3x+1$ problem: new lower bounds on nontrivial cycle lengths,''
\textit{Discrete Math.}\ \textbf{118} (1993), 45--56.

\bibitem{Hales17}
T.~Hales et al., ``A formal proof of the Kepler conjecture,''
\textit{Forum Math.\ Pi} \textbf{5} (2017), e2.

\bibitem{Hercher25}
J.\,C.~Hercher, ``There are no Collatz $m$-cycles with $m \leq 91$,''
\textit{J.~Integer Seq.}\ \textbf{28} (2025), Article 25.2.2.

\bibitem{Lagarias85}
J.\,C.~Lagarias, ``The $3x+1$ problem and its generalizations,''
\textit{Amer.\ Math.\ Monthly} \textbf{92} (1985), 3--23.

\bibitem{LMN95}
M.~Laurent, M.~Mignotte, and Yu.~Nesterenko, ``Formes lin\'eaires en deux logarithmes et d\'eterminants d'interpolation,''
\textit{J.~Number Theory} \textbf{55} (1995), 285--321.

\bibitem{Matveev00}
E.\,M.~Matveev, ``An explicit lower bound for a homogeneous rational linear form in the logarithms of algebraic numbers, II,''
\textit{Izv.\ Math.}\ \textbf{64} (2000), 1217--1269.

\bibitem{SdW05}
D.~Simons and B.~de~Weger, ``Theoretical and computational bounds for $m$-cycles of the $3n+1$ problem,''
\textit{Acta Arith.}\ \textbf{117} (2005), 51--70.

\bibitem{Steiner77}
R.\,P.~Steiner, ``A theorem on the Syracuse problem,''
in \textit{Proc.\ 7th Manitoba Conf.\ Numer.\ Math.\ Comput.}, 1977, pp.~553--559.

\bibitem{Tao22}
T.~Tao, ``Almost all orbits of the Collatz map attain almost bounded values,''
\textit{Forum Math.\ Pi} \textbf{10} (2022), e12.

\bibitem{Terras76}
R.~Terras, ``A stopping time problem on the positive integers,''
\textit{Acta Arith.}\ \textbf{30} (1976), 241--252.

\bibitem{Wielandt50}
H.~Wielandt, ``Unzerlegbare, nicht negative Matrizen,''
\textit{Math.\ Z.}\ \textbf{52} (1950), 642--648.

\end{thebibliography}

\end{document}
